% © 2026 Ing. David Jaroš — CC BY-NC-ND 4.0
%
% This work is licensed under a Creative Commons Attribution-NonCommercial-NoDerivatives
% 4.0 International License (CC BY-NC-ND 4.0).
%
% License History: Earlier drafts (up to v0.3) were released under CC BY 4.0.
% From v0.4 onward, all material is released under CC BY-NC-ND 4.0 to protect
% the integrity of the theoretical work during ongoing academic development.
%
% See LICENSE.md for full license text.

\documentclass[12pt,a4paper]{article}
\usepackage[a4paper,margin=2.5cm]{geometry}
\usepackage{amsmath,amssymb,amsthm}
\usepackage{booktabs}
\usepackage{tcolorbox}
\tcbuselibrary{skins,breakable}

\newtheorem{proposition}{Proposition}[section]
\newtheorem{remark}[proposition]{Remark}
\newtcolorbox{statusbox}[1][]{colback=gray!8!white,colframe=black!60,arc=3pt,boxrule=1pt,title=Status Summary,#1}

\title{Fermionic Sector from Gauge-Consistency Constraints}
\author{Ing.\ David Jaroš}
\date{2026-05-12}

\begin{document}
\maketitle

\section{Goal}
This note addresses the alternate route to fermions:
derive constraints from gauge consistency/anomaly cancellation, rather than
factorizing the bosonic kinetic action.

\section{Step 1: anomaly constraints from gauge consistency}
For a chiral gauge theory, consistency requires vanishing anomaly polynomials.
For a U(1) charge assignment \(q_f\), the cubic and linear conditions are
\[
\sum_f q_f^3=0,\qquad \sum_f q_f=0.
\]
In the electroweak sector this is applied to hypercharge assignments of
left-handed Weyl fields (including conjugate right-handed singlets).

\begin{proposition}[UBT-consistent SM-like hypercharge set is anomaly-free]
Using one-generation chiral content
\[
Q_L,\ u_R^c,\ d_R^c,\ L_L,\ e_R^c
\]
with
\[
Y=\Bigl(\tfrac16,\ -\tfrac23,\ \tfrac13,\ -\tfrac12,\ 1\Bigr),
\]
including colour/SU(2) multiplicities, both
\(
\sum_f Y_f^3
\)
and
\(
\sum_f Y_f
\)
vanish.
\end{proposition}

\begin{proof}
Multiplicity-weighted sums per generation:
\[
\sum_f Y_f
=6\!\cdot\!\tfrac16 + 3\!\cdot\!\Bigl(-\tfrac23\Bigr) + 3\!\cdot\!\tfrac13
+2\!\cdot\!\Bigl(-\tfrac12\Bigr) + 1\!\cdot\!1 = 0,
\]
\[
\sum_f Y_f^3
=6\!\cdot\!\Bigl(\tfrac16\Bigr)^3
+3\!\cdot\!\Bigl(-\tfrac23\Bigr)^3
+3\!\cdot\!\Bigl(\tfrac13\Bigr)^3
+2\!\cdot\!\Bigl(-\tfrac12\Bigr)^3
+1^3 = 0.
\]
\end{proof}

\section{Step 2: mixed gravitational/electroweak check}
For the mixed \([\mathrm{SU}(2)_L]^2U(1)_Y\) condition:
\[
\sum_{\text{SU(2) doublets}} T(r_f)\,Y_f = 0,
\quad T(\mathbf{2})=\tfrac12.
\]
With three coloured quark doublets and one lepton doublet:
\[
3\cdot \tfrac12\cdot \tfrac16 + 1\cdot \tfrac12\cdot\Bigl(-\tfrac12\Bigr)=0.
\]
So the mixed condition is satisfied for the same chiral set.

\section{Step 3: WZW interpretation status}
The WZW counterterm route is a candidate mechanism for encoding anomaly
matching in an effective bosonic action. In current UBT files, this remains
model-construction level [MC]: existence of a fully derived
\(
S_{\mathrm{WZW}}[\Theta]
\)
compatible with \(\CC\otimes\HH\) is not yet closed at [L1].

\begin{remark}
This route derives \emph{consistency constraints} on fermionic content.
It does not yet provide a full first-principles derivation of the fermionic
kinetic term from canonical \(S[\Theta]\).
\end{remark}

\begin{statusbox}
Gauge anomaly-cancellation in UBT: \textbf{[satisfied for SM-like chiral content; L1 check]}\\
Fermionic content from anomaly-cancellation: \textbf{[partially derived as a consistency requirement]}\\
WZW člen v $S[\Theta]$: \textbf{[OPEN/MC]}\\
$N_{\mathrm{helicity}}=2$ z gaugeové konzistence: \textbf{[MC]}\\
Dopad: $N_{\mathrm{eff}}=12$ a $B_0=8\pi$ \textbf{[zůstává OPEN]}.
\end{statusbox}

\end{document}
